% !TeX root = ../../Book2.tex
\section{Bochner–Riesz 问题}
\label{b2:sec:2403}

可积核条件给出同时适用于所有 \(p\) 的范围，却未必给出固定 \(p\) 所需的最小平滑量。中心指数 \(p=2\) 只要求符号有界；离开中心后，核的尾部和不同方向的集中分别产生障碍。本节先给出可以直接检验的必要条件，再证明一个解析插值充分范围，最后在已知乘子有界性的前提下量化逼近速度。必要条件、充分条件和逼近阶数承担不同职责，不能用其中一个代替另外两个。

\subsection{\texorpdfstring{\(L^2\)}{L2} 理论}
\label{b2:subsec:240301}

设 \(d\geq1\)，固定 \(\delta\geq0\)。对每个 \(R>0\)，\cref{b2:def:24-bochner-riesz}中的符号
\(m_{\delta,R}(\xi)=(1-|\xi|^2/R^2)_+^\delta\)
在球内介于零和一之间，在原点邻域任意接近一；零阶时它是球的示性函数。因此符号的本性上确界恰为 \(1\)，\cref{b2:prop:22-l2}给出
\[
 \|B_R^\delta\|_{2\to2}=1,\quad
 \lim_{R\to\infty}\|B_R^\delta f-f\|_2=0
 \quad(f\in L^2).
\]
收敛的具体理由是 \(m_{\delta,R}(\xi)\to1\)，而
\(|m_{\delta,R}-1|^2\cdot|\widehat f|^2\leq|\widehat f|^2\in L^1\)；
先在 Fourier 一侧用支配收敛，再用 Plancherel 返回空间侧。
但 \(L^2\) 范数界不控制频率边界上的点态取值，也不自动控制 \(L^p\)、\(p\ne2\)。与限制问题的联系需要额外结构。

\ProseParagraphBreak
对于一维 \(\delta=0\)，球就是区间，Hilbert 变换给出所有 \(1<p<\infty\) 上的界。高维球不再是有限个半轴投影的乘积。矩形投影可以逐坐标应用 Hilbert 变换并用 Fubini 得到有界性；球边界的法向连续变化，出现了无限方向的组合。

\ProseParagraphBreak
平方理论还允许量化平均逼近的速度，而且显示有无边界消失阶会改变速度结构。
\Needspace{6\baselineskip}
\begin{proposition}\label{b2:prop:24-l2-rate}
设整数 \(d\geq1\)，\(B_R^\delta\) 是 \(L^2(\mathbb R^d;\mathbb C)\) 上符号为 \((1-|\xi|^2/R^2)_+^\delta\) 的 Fourier 乘子，\(R>0\)；零阶符号取开球的示性函数。设 \(\delta>0\)、\(0<s\leq2\)，且
\(f\in L^2(\mathbb R^d;\mathbb C)\) 满足
\((2\pi|\xi|)^s\cdot\widehat f(\xi)\in L^2\)。
以这一符号定义 \((-\Delta)^{s/2}f\)，则
\[
 \|B_R^\delta f-f\|_2
 \leq C_\delta(2\pi R)^{-s}\cdot\|(-\Delta)^{s/2}f\|_2.
\]
若 \(\delta=0\)，则对任意 \(s>0\) 同一结论成立，常数可以取一。
\end{proposition}
\begin{proof}
对 \(a\geq0\)，初等均值估计给出
\[
 |1-(1-a)_+^\delta|\leq\max\{1,\delta\}\min\{1,a\}.
\]
对 \(0<\delta<1\)，在 \(0\leq a\leq1\) 上有
\((1-a)^\delta\geq1-a\)，所以误差不超过 \(a\)；对 \(\delta\geq1\)，同一区间上的导数上界为 \(\delta\)。当 \(a\geq1\)，误差等于一。令 \(a=|\xi|^2/R^2\)。因为 \(0<s\leq2\)，有
\(\min\{1,a\}\leq a^{s/2}\)，所以
\[
 |m_{\delta,R}(\xi)-1|
 \leq C_\delta R^{-s}|\xi|^s.
\]
将两边乘以 \(|\widehat f|\)，取平方积分并用 Plancherel，即得第一项；因 \((-\Delta)^{s/2}\) 的符号是 \((2\pi|\xi|)^s\)，常数中出现 \((2\pi R)^{-s}\)。零阶硬截断的误差只在 \(|\xi|>R\)，在那里
\(\mathbf1_{\{|\xi|>R\}}\leq(|\xi|/R)^s\)
对任意 \(s>0\) 成立，故得到第二项。
\end{proof}
对于 \(\delta>0\)，在一个固定紧频带内有
\(R^2[1-(1-|\xi|^2/R^2)^\delta]\to\delta|\xi|^2\)。
若非零 \(\widehat f\) 支于该频带，支配收敛于是给出
\[
 R^2(B_R^\delta f-f)\longrightarrow
 -\frac{\delta}{4\pi^2}(-\Delta)f
 \quad\text{于 }L^2.
\]
因此通常不能仅靠输入具有更多光滑性把第一项误差提高到 \(R^{-s}\)、\(s>2\)。要提高阶数，应组合不同平均以消去 Taylor 首项，或者采用在低频附近恰好等于一的另一种截断。频率边界的平滑与低频的逼近阶数是两个不同因素。

\ProseParagraphBreak

这个例子也说明“越平滑，逼近必然越快”不是无条件原则。硬截断在固定紧频带上最终误差为零；正阶平均虽然在球边界更光滑，却在球内部偏离一，因而留下低阶 Taylor 误差。硬截断的困难主要是其他 \(L^p\) 的一致有界性和逐点求和；平方逼近本身则十分直接。

\subsection{临界指标}
\label{b2:subsec:240302}

\Needspace{6\baselineskip}
\begin{proposition}[尾部必要条件]\label{b2:prop:24-necessary}
设整数 \(d\geq1\)、\(\delta\geq0\)、\(1\leq p\leq\infty\)，\(B_1^\delta\) 为符号 \((1-|\xi|^2)_+^\delta\) 的 \(L^2\) Fourier 乘子，零阶符号取单位开球的示性函数。若它具有与 Schwartz 输入上的 Fourier 定义一致的有界延拓 \(L^p(\mathbb R^d;\mathbb C)\to L^p(\mathbb R^d;\mathbb C)\)，则在 \(p\ne2\) 时必须满足
\[
 \delta>d\left|\frac1p-\frac12\right|-\frac12.
\]
当右侧为负数时，该条件不提供比 \(\delta\ge0\) 更强的限制。
\end{proposition}
\begin{proof}
选 \(\widehat f\in C_c^\infty\)，在闭单位球附近等于 \(1\)。则 \(f\) 为 Schwartz 函数，且
\[
 B_1^\delta f=K_\delta
 =\mathcal F^{-1}\bigl[(1-|\xi|^2)_+^\delta\bigr].
\]
所以对 \(p<\infty\)，算子有界性迫使 \(K_\delta\in L^p\)。
\cref{b2:prop:24-kernel-tail}给出的精确核可积性条件为
\[
 p\bigl(\delta+(d+1)/2\bigr)>d,
 \quad\text{即}\quad
 \delta>\frac dp-\frac{d+1}{2}.
\]
恢复严格不等式的来源：核主项的余弦在每个足够大的单位径向区间内都有一段固定比例集合保持模下界，余项比主项低一阶。故径向绝对积分至少为下列级数的某个正常数倍：
\[
 \sum_{n\geq n_0}n^{d-1-p(\delta+(d+1)/2)}.
\]
阈值等号时正是发散的调和级数。

符号为实值，在 Schwartz 类上有
\(\langle B_1^\delta f,g\rangle=\langle f,B_1^\delta g\rangle\)。
若 \(1<p<\infty\)，由此和 \(L^p\) 对偶性，\(p\) 上的界也给出 \(p'\) 上相同的界。分别应用刚得的尾部条件，再取较强者，得到
\[
 \delta>\max\left\{\frac dp,\frac d{p'}\right\}
 -\frac{d+1}{2}
 =d\left|\frac1p-\frac12\right|-\frac12.
\]

若 \(p=1\)，直接得到 \(\delta>(d-1)/2\)。若 \(p=\infty\)，对 Schwartz 函数 \(f,g\) 使用实符号配对和假设的 \(L^\infty\) 界，有
\[
 |\langle B_1^\delta f,g\rangle|
 =|\langle f,B_1^\delta g\rangle|
 \leq C\|f\|_1\cdot\|g\|_\infty.
\]
对紧支撑光滑 \(g\) 的单位球取上确界，便得到
\(\|B_1^\delta f\|_1\leq C\|f\|_1\)，于是也可应用 \(p=1\) 的必要条件。给定延拓在 Schwartz 类上的配对关系与 \(L^1\) 范数的对偶公式，便把无穷指数的问题转到可积端点。对于 \(p=2\)，尾部检验给出 \(\delta>-1/2\)，所有 \(\delta\geq0\) 都满足。
\end{proof}
严格性来自径向调和级数；没有对尾部作精确下界，就不能宣称阈值等号也失败。这个检验并未排除临界函数
\(\max\{d|1/p-1/2|-1/2,0\}\) 中为零的内部指数上的硬截断；下一小节的深层几何定理才排除高维硬球。

\subsection{Bochner–Riesz 猜想}
\label{b2:subsec:240303}

经典 Bochner–Riesz 目标是：对 \(d\ge2\)、\(1<p<\infty\)，在
\[
 \delta>\max\left\{d\left|\frac1p-\frac12\right|-\frac12,\ 0\right\}
\]
时取得 \(L^p\) 有界性。这里严格的正平滑量将硬球情形排除在目标外。维数、端点形式、径向输入限制以及局部版本各是不同问题。

\ProseParagraphBreak

下面采用\term[fefferman-qiuchengzi-dingli]{Fefferman 球乘子定理}[Fefferman ball multiplier theorem]的下述版本。

\Needspace{6\baselineskip}
\begin{theorem}[深层结构定理：Fefferman 球乘子定理]\label{b2:thm:24-fefferman}
设整数 \(d\geq2\)，\(1<p<\infty\)。在 \(L^2(\mathbb R^d;\mathbb C)\) 上以符号
\(\mathbf1_{\{\,|\xi|<1\,\}}\) 定义单位球 Fourier 乘子。它具有与 Schwartz 类上一致的有界延拓
\(L^p(\mathbb R^d;\mathbb C)\to L^p(\mathbb R^d;\mathbb C)\)，当且仅当 \(p=2\)。
\end{theorem}
\begin{proofsketch}
假设球乘子在某个非二指数有界。对球作放大、频移并取极限，配合随机符号，可推出各方向半空间投影的一致向量值界。Besicovitch 型细矩形构造使输入的平方函数质量分散，而输出在一个小测度集合中累积；对 \(p>2\) 这与向量值界矛盾。对偶性处理 \(p<2\)，切片传递处理高维。
几何矩形构造及传递论证的详细证明见原文
\cite[Theorem 1; Section 2, pp.~330--336]{Fefferman1971Ball}。
\end{proofsketch}
这个定理并不说任何正的 \(\delta\) 都足够。球乘子的失败与正平滑后的范围之间，仍隔着一个依赖维数与 \(p\) 的估计问题。

\subsection{解析插值与 \texorpdfstring{\(L^p\)}{Lᵖ} 逼近估计}
\label{b2:subsec:240304}

可积核给出了 \(L^1\) 和 \(L^\infty\) 两个强端点，平方理论给出了中心指数。它们对应不同的平滑阶数；要连接它们，必须同时改变算子中的参数，而不能只对一个未知有界的固定算子插值。先证明复参数竖线上的核估计，再把参数与输入指数联系起来。

\Needspace{6\baselineskip}
\begin{lemma}[固定实部的复阶 Bochner--Riesz 核]\label{b2:lem:24-complex-kernel}
设整数 \(d\geq1\)，固定实数 \(\sigma>-1\)。对 \(z=\sigma+\mathrm it\)、\(t\in\mathbb R\)，令
\[
 m_z(\xi)=
 \begin{cases}
  \exp\bigl(z\log(1-|\xi|^2)\bigr),&|\xi|<1,\\
  0,&|\xi|\geq1,
 \end{cases}
 \quad K_z=\mathcal F^{-1}m_z,
\]
其中球内使用实对数。符号 \(m_z\) 属于 \(L^1\)，所以 \(K_z\) 是有界连续函数。存在只依赖 \(d,\sigma\) 的常数 \(C\)，使
\[
 |K_{\sigma+\mathrm it}(x)|
 \leq C e^{C(1+t^2)}
 (1+|x|)^{-\sigma-(d+1)/2},
 \quad x\in\mathbb R^d,\ t\in\mathbb R.
\]
特别地，若 \(\sigma>(d-1)/2\)，则
\(\|K_{\sigma+\mathrm it}\|_1\leq C e^{C(1+t^2)}\)。
\end{lemma}
\begin{proof}
记 \(\nu=d/2+z=u+\mathrm it\)，其中 \(u=d/2+\sigma>-1/2\) 固定。\subsecref{b2:subsec:230402}已经以整函数 \(P=1/\Gamma\) 构造复阶 Bessel 级数，并证明归一化球内分布的核恒等式。这里 \(\operatorname{Re}z>-1\)，是普通可积符号，乘回 Gamma 因子即得
\[
 K_z(x)=\Gamma(z+1)\pi^{-z}
 |x|^{-\nu}J_\nu(2\pi|x|),\quad x\ne0.
\]
这是未归一化平均的核，因子 \(\Gamma(z+1)\) 必须保留。复 Gamma 的积分定义给出
\[
 |\Gamma(\sigma+1+\mathrm it)|
 \leq\int_0^\infty e^{-v}v^\sigma\,\mathrm dv
 =\Gamma(\sigma+1).
\]
因此分子 Gamma 因子在这条竖线上甚至一致有界。

小半径与初值由级数控制。对 \(0<r\leq1\)，
\[
 r^{-\nu}J_\nu(2\pi r)
 =\pi^\nu\sum_{k=0}^\infty
 \frac{(-\pi^2r^2)^k}{k!\Gamma(k+\nu+1)}.
\]
因为 \(u+1>0\)，Gamma 递推给出
\[
 |\Gamma(k+\nu+1)|^{-1}
 \leq
 \frac{|P(u+1+\mathrm it)|}
 {\prod_{\ell=0}^{k-1}(u+1+\ell)}.
\]
分母是增长的正实乘积；配合 \(k!\)，级数及其任意固定阶 \(r\) 导数均可在 \(0\leq r\leq1\) 一致求和。前章从 Gamma 无限乘积得到
\[
 |P(u+1+\mathrm it)|\leq C_u e^{C_u(1+t^2)}.
\]
于是 \(|r^{-\nu}J_\nu(2\pi r)|\leq C e^{C(1+t^2)}\)，并且
\[
 |J_\nu(2\pi)|+
 \left|\frac{\mathrm d}{\mathrm dr}J_\nu(2\pi r)
 \bigg|_{r=1}\right|
 \leq C e^{C(1+t^2)}.
\]
对 \(r^\nu\) 求导带来的 \(1+|\nu|\) 因子也可并入此界。前面的 Gamma 分子界因而给出 \(0\leq|x|\leq1\) 的目标估计，原点按连续极限解释。

对 \(r\geq1\)，令
\[
 w(r)=r^{1/2}J_\nu(2\pi r),\quad
 E(r)=|w'(r)|^2+4\pi^2|w(r)|^2.
\]
复阶 Bessel 方程仍给出
\[
 w''+4\pi^2w=(\nu^2-1/4)r^{-2}w.
\]
能量 \(E\) 是实的非负量，不需要系数 \(\nu^2\) 为实数。微分并消去自由振荡项后，
\[
 |E'(r)|
 \leq C|\nu^2-1/4|r^{-2}\cdot E(r)
 \leq C_u(1+t^2)r^{-2}\cdot E(r).
\]
积分 \(\int_1^\infty r^{-2}\,\mathrm dr=1\)，结合已证的 \(r=1\) 初值界，得到
\[
 E(r)\leq E(1)e^{C_u(1+t^2)}
 \leq C e^{C(1+t^2)},\quad r\geq1.
\]
因此
\[
 |J_\nu(2\pi r)|
 \leq C e^{C(1+t^2)}r^{-1/2}.
\]
核公式中另外的幂因子模为 \(\pi^{-\sigma}r^{-u}\)，所以大半径衰减为
\(r^{-u-1/2}=r^{-\sigma-(d+1)/2}\)。
与小半径估计合并即得结论。最后，当
\(\sigma+(d+1)/2>d\) 时，径向积分可积，便得到统一的 \(L^1\) 核界。
\end{proof}

与球面限制中的竖线不同，此处 \(\nu\) 的实部是 \(d/2+\sigma\)。能量法仍然适用，因为对固定实部有
\(|\nu^2-1/4|\leq C(1+t^2)\)；不能直接把只对固定实阶陈述的 Bessel 渐近当作这一统一估计。另一个区别在插值的目标：本节要得到 \(L^p\to L^p\)，而\cref{b2:lem:23-compact-analytic-interpolation}的陈述给出的是 \(L^p\to L^{p'}\)。下面保留其共同定义域的做法，重新选择测试函数指数。

\Needspace{6\baselineskip}
\begin{theorem}[可直接证明的插值范围]\label{b2:thm:24-interpolation-range}
设整数 \(d\geq1\)、\(1<p<\infty\)，实数 \(\delta\) 满足
\[
 \delta>(d-1)\left|\frac1p-\frac12\right|,
\]
并在 Schwartz 类上以
\[
 B_R^\delta f=\mathcal F^{-1}
 \bigl[(1-|\xi|^2/R^2)_+^\delta\cdot\widehat f(\xi)\bigr],
 \quad R>0,
\]
定义 Bochner--Riesz 平均。则这些算子具有与 \(L^2\) 定义相容的唯一的 \(L^p(\mathbb R^d;\mathbb C)\) 有界延拓，且
\[
 \sup_{R>0}\|B_R^\delta\|_{p\to p}\leq C_{d,p,\delta},
 \quad
 \lim_{R\to\infty}\|B_R^\delta f-f\|_p=0
 \quad(f\in L^p).
\]
\end{theorem}
\begin{proof}
当 \(p=2\) 时，结论已由本节开头的 Plancherel 论证得到，甚至允许 \(\delta=0\)。先设 \(1<p<2\)，令
\[
 \theta=\frac2p-1\in(0,1),\quad
 \sigma=\frac\delta\theta.
\]
题设恰好使 \(\sigma>(d-1)/2\)，而目标参数满足 \(\delta=\theta\sigma\)。在闭条带 \(0\leq\operatorname{Re}w\leq1\) 上，考察
\[
 S_wf=\mathcal F^{-1}(m_{\sigma w}\cdot\widehat f),
 \quad
 m_{\sigma w}(\xi)=(1-|\xi|^2)_+^{\sigma w}.
\]
球内使用实对数，球外为零，球面零测集的取值无关。在整条带上
\(|m_{\sigma w}|\leq1\)，且符号支撑于闭单位球。
对内点参数的任意固定阶导数，只多出
\(\sigma^k\log^k(1-|\xi|^2)\)；
在参数紧集上，其绝对积分由球面法向变量中可积的
\(u^\varepsilon|\log u|^k\) 控制。因此符号与光滑频率测试函数的配对全纯，并连续到两条边界。

先明确一套可随参数变化的共同测试空间。令 \(f,g\) 是有界紧支撑简单函数。它们的 Fourier 变换有界，所以
\[
 B_w(f,g)=\int_{|\xi|<1}
 m_{\sigma w}(\xi)\cdot\widehat f(\xi)
 \cdot\widehat g(-\xi)\,\mathrm d\xi
\]
是绝对收敛的双线性配对。它也等于
\(\int(S_wf)(x)\cdot g(x)\,\mathrm dx\)：符号可积使核有界，两个输入紧支撑，故两种写法由 Fubini 相等。这个共同定义域包括稍后使用的所有变化测试函数。

取 \(\|f\|_p=\|g\|_{p'}=1\)，把它们的非零值分别写成
\(f=\sum_j a_j\mathbf1_{E_j}\)、
\(g=\sum_k b_k\mathbf1_{F_k}\)，同一族内的集合两两不交且有界。定义
\[
 \begin{aligned}
 f_w&=\sum_j\frac{a_j}{|a_j|}
 |a_j|^{p(1+w)/2}\cdot\mathbf1_{E_j},\\
 g_w&=\sum_k\frac{b_k}{|b_k|}
 |b_k|^{p'(1-w)/2}\cdot\mathbf1_{F_k}.
 \end{aligned}
\]
在 \(w=\theta\) 时，两个复幂指数都等于一，所以
\(f_\theta=f\)、\(g_\theta=g\)。在竖边上则有
\[
 \|f_{\mathrm it}\|_2=\|g_{\mathrm it}\|_2=1,\quad
 \|f_{1+\mathrm it}\|_1=1,\quad
 \|g_{1+\mathrm it}\|_\infty=1.
\]
测试函数 \(g_w\) 的指数与 \(f_w\) 方向相反，正是为了在右边界配对 \(L^1\) 输出与 \(L^\infty\) 测试。

函数 \(F(w)=B_w(f_w,g_w)\) 在条带内全纯并连续到边界，因为 \(f_w,g_w\) 都是有限解析组合。并且各系数的模仅依赖 \(\operatorname{Re}w\)，于是
\[
 |F(w)|\leq |B(0,1)|\cdot
 \|f_w\|_1\cdot\|g_w\|_1\leq C_{f,g}
\]
在整个闭条带成立。有限解析组合及这项统一上界共同保证了条带插值所需的解析性与增长控制。

在 \(w=\mathrm it\) 上，符号模不超过一，Plancherel 与 Cauchy--Schwarz 给出
\(|F(\mathrm it)|\leq1\)。
在 \(w=1+\mathrm it\) 上，\cref{b2:lem:24-complex-kernel}及 Young 不等式给出
\[
 \|S_{1+\mathrm it}\|_{1\to1}
 \leq C e^{C(1+\sigma^2t^2)},
 \quad
 |F(1+\mathrm it)|\leq C e^{C(1+\sigma^2t^2)}.
\]
取固定 \(\kappa>C\sigma^2\)，对
\(e^{\kappa(w-\theta)^2}F(w)\)
应用\cref{b2:thm:17-three-lines}。这个因子在竖直方向压住右边界增长，在插值点等于一，边界的实部只带来固定常数。由此得到
\[
 |B_\theta(f,g)|\leq C_{d,p,\delta}.
\]
齐次性与 \(L^p\) 对偶性说明
\(\|S_\theta f\|_p\leq C_{d,p,\delta}\|f\|_p\)
在紧支撑简单函数上成立；稠密性给出唯一的 \(L^p\) 延拓。
它与原 Fourier 定义相容：对 Schwartz 输入作有界紧支撑简单函数逼近，可同时在 \(L^1,L^2,L^p\) 中收敛，频率配对与 \(L^2\) 乘子极限均保持一致。因此 \(S_\theta=B_1^\delta\)。

若 \(2<p<\infty\)，则 \(1<p'<2\)，并且
\(|1/p'-1/2|=|1/p-1/2|\)。
实值符号 \(m_\delta\) 使 \(B_1^\delta\) 在测试类上自伴随，已证的 \(L^{p'}\) 界通过对偶性给出 \(L^p\) 界。最后，\cref{b2:prop:24-dilation}的同指数共轭公式使各半径范数相同，\cref{b2:prop:24-strong-approximation}用统一算子界与稠密测试类给出 \(R\to\infty\) 时的强收敛。
\end{proof}
这是一个明确已证范围，一般窄于经典目标。它的证明只结合 \(L^2\) 与可积核端点，没有利用球面限制的额外改进。比较两条临界直线，能看清深层理论应在哪个区间增加信息。

\ProseParagraphBreak
已经证明的一致乘子界也能与输入正则性配合，给出其他 \(L^p\) 中的逼近速度。
\Needspace{6\baselineskip}
\begin{proposition}\label{b2:prop:24-lp-rate}
设整数 \(d\geq1\)，\(1<p<\infty\)，\(\delta>0\)，\(0<s\leq2\)。设符号 \((1-|\xi|^2/R^2)_+^\delta\) 的 Fourier 乘子具有与 \(L^2\) 定义相容的延拓 \(B_R^\delta:L^p\to L^p\)，并满足
\[
 M\coloneqq\sup_{R>0}\|B_R^\delta\|_{p\to p}<\infty.
\]
令 \(H^{s,p}(\mathbb R^d)\) 按\cref{b2:def:22-bessel-space}定义，即
\(\|f\|_{H^{s,p}}=\|J^sf\|_p\)，其中 \(J^s\) 的符号为
\((1+4\pi^2|\xi|^2)^{s/2}\)。则对每个 \(f\in H^{s,p}(\mathbb R^d)\) 和 \(R\geq1\)，
\[
 \|B_R^\delta f-f\|_p
 \leq C_{d,p,s,\delta,M}R^{-s}\cdot\|f\|_{H^{s,p}}.
\]
\end{proposition}
\begin{proof}
先对 Schwartz 输入证明；该类在 \(H^{s,p}\) 中稠密，且 \(J^{-s}\) 在 \(L^p\) 上有界，最终可把估计传到一般输入。取固定径向
\(\chi\in C_c^\infty(\mathbb R^d)\)，在 \(|\xi|\leq1/4\) 上为一，支撑于 \(|\xi|<1/2\)。令
\(P_Rf=\mathcal F^{-1}[\chi(\xi/R)\cdot\widehat f]\)，并记 \(g=J^sf\)。分解
\[
 (I-B_R^\delta)f
 =(I-B_R^\delta)P_Rf+
 (I-B_R^\delta)(I-P_R)f.
\]
后一个输入只含 \(|\xi|\geq R/4\) 的高频。把其 \(R^{-s}\) 因子提出后，作用于 \(g\) 的乘子为
\[
 a_R(\xi)=R^s(1-\chi(\xi/R))
 (1+4\pi^2|\xi|^2)^{-s/2}.
\]
在其支撑上 \(R\leq4|\xi|\)。每次微分 Bessel 因子带来相应的 \((1+|\xi|)^{-1}\)，微分截断的项只出现在 \(|\xi|\asymp R\) 上。因此对 \(|\alpha|\leq d+2\)，
\[
 |\xi|^{|\alpha|}\cdot|\partial^\alpha a_R(\xi)|
 \leq C_\alpha R^s(1+|\xi|)^{-s}\leq C_\alpha
\]
一致成立。\cref{b2:thm:22-mihlin}给出
\(\|(I-P_R)f\|_p\leq C R^{-s}\|g\|_p\)，再应用
\(\|I-B_R^\delta\|_{p\to p}\leq1+M\)
便控制高频误差。

低频误差的归一化乘子为
\[
 \ell_R(\xi)=R^s[1-(1-|\xi|^2/R^2)^\delta]\cdot
 \chi(\xi/R)(1+4\pi^2|\xi|^2)^{-s/2}.
\]
在截断支撑上 \(v=|\xi|^2/R^2\leq1/4\)，远离符号的球边界。由微积分基本定理，
\[
 1-(1-v)^\delta=v\,b_\delta(v),\quad
 b_\delta(v)=\delta\int_0^1(1-tv)^{\delta-1}\,\mathrm dt.
\]
所以 \(b_\delta\) 及每个所需阶数的导数在 \(0\leq v\leq1/4\) 有界。
把低频乘子写成
\[
 \ell_R(\xi)=R^{s-2}|\xi|^2\cdot
 b_\delta(|\xi|^2/R^2)\cdot\chi(\xi/R)
 (1+4\pi^2|\xi|^2)^{-s/2}.
\]
链式与乘积求导可分成两个区域记录：对所有 \(|\alpha|\leq d+2\)，
\[
 |\xi|^{|\alpha|}\cdot|\partial^\alpha\ell_R(\xi)|
 \leq C_\alpha
 \begin{cases}
 R^{s-2}|\xi|^2,&0<|\xi|\leq1,\\
 (|\xi|/R)^{2-s},&1\leq|\xi|\leq R/2.
 \end{cases}
\]
第一行利用了 \(|\xi|^2\) 的零点和其余因子在单位球内的光滑有界性；高阶导数即使消去这个零点，外面的 \(|\xi|^{|\alpha|}\) 也补足所需幂次。第二行中，微分 \(b_\delta\) 或截断产生的尺度因子由 \(R^{-1}\lesssim|\xi|^{-1}\) 控制，微分 Bessel 因子带来 \(|\xi|^{-1}\)。这解释了每次求导后仍保持同一零阶界的原因。

因为 \(R\geq1\)、\(s\leq2\)，两行都一致有界；支撑外乘子为零。再用 Mihlin 定理，得到低频误差不超过 \(C R^{-s}\|g\|_p\)。合并两部分后，以 \(H^{s,p}\) 中的 Schwartz 逼近传到一般 \(f\)，即得结论。
\end{proof}
在一致乘子界成立的指数范围内，这个论证把两种误差分开：高频尾部由输入的光滑性控制，低频误差由平均符号在原点的 Taylor 阶数控制。若进一步构造消掉低阶 Taylor 系数的组合平均，后一限制可以改变，但需要重新核对组合的统一算子界。

\ProseParagraphBreak

例如，对固定 \(d,p\) 选择处于已证插值范围的 \(\delta\)，由\cref{b2:prop:22-integer-sobolev}的整数阶范数等价，一个 \(W^{1,p}\) 输入有 \(R^{-1}\) 的范数逼近速度，一个 \(W^{2,p}\) 输入有 \(R^{-2}\) 的速度。这里只是范数误差，不能把它直接变成全参数逐点误差率；后者需要位置上的极大控制或更强的正则性与嵌入。

\ProseParagraphBreak

必要直线与插值直线的距离可以用一个数值模型说明。在 \(d=3\)、\(p=4\) 中，尾部条件要求 \(\delta>1/4\)，本章直接证明的插值范围要求 \(\delta>1/2\)。例如 \(\delta=3/4\) 在已证范围内，\(\delta=1/8\) 被必要条件排除，而 \(\delta=3/8\) 不能仅靠这两项结果判断。说它“没有被反例排除”与说它“已证明有界”有明确区别。

\ProseParagraphBreak

指数对偶把 \(p=4\) 与 \(p=4/3\) 的范围对应起来，因为球符号为实且算子在测试类上自伴随。缩放又使每个已证同指数界自动对所有半径具有同一常数。因此进一步扩大范围的关键，是获得更精细的平滑量与输入指数关系。若一篇文献只给固定空间球内的局部输出或径向输入，其命题还不能直接填入这个全空间、一般输入的空缺。

\ProseParagraphBreak

在 \(p\) 接近二时，尾部必要条件的右侧可能为负，此时它只要求原来已经约定的 \(\delta\geq0\)，看起来没有门槛。然而高维 \(\delta=0\) 已由球乘子定理排除所有非二指数。这个额外的几何障碍来自多方向叠加，并不是单个核尾部的一条更仔细积分能看出来的。正阶问题是否能在接近二的范围内降到任意小的平滑量，则要依赖更深的曲率和限制估计。

\ProseParagraphBreak

对角线上的 \(p=2\) 则没有这些困难：符号有界已经给出范数一，且所有 \(\delta\geq0\) 都随固定阶数的半径增加强收敛到恒等。这个特殊点是各种插值证明的可靠起点，但它不控制其他指数上的几何组合。把一个证明完全归结为“平方理论加插值”时，应能够说清另一个端点来自哪里，以及为什么它有足够强的范围和统一常数。
\ProseParagraphBreak
消去低频误差的组合可以直接构造，而无需增加单个平均的边界消失阶。半径放大到 \(\sqrt2R\) 时，低频二次项的系数减半；用权 \(2,-1\) 相减，就能保留常数项并消去该误差。

\Needspace{6\baselineskip}
\begin{proposition}[两半径组合的逼近阶数]\label{b2:prop:24-two-radius-rate}
设整数 \(d\geq1\)，\(1<p<\infty\)，\(\delta>0\)，且符号
\((1-|\xi|^2/R^2)_+^\delta\) 的相容 \(L^p\) Fourier 乘子满足
\(\sup_{R>0}\|B_R^\delta\|_{p\to p}\leq M<\infty\)。
定义
\[
 C_R=2B_{\sqrt2R}^\delta-B_R^\delta,\quad R>0.
\]
则 \(\|C_R\|_{p\to p}\leq3M\)，并且对 \(0<s\leq4\)、
\(R\geq1\) 和 \(f\in H^{s,p}(\mathbb R^d)\)，有
\[
 \|C_Rf-f\|_p
 \leq C_{d,p,s,\delta,M}R^{-s}\cdot\|f\|_{H^{s,p}}.
\]
这里 \(H^{s,p}\) 及其范数按\cref{b2:def:22-bessel-space}定义。
若 \(\delta=1\)，上述速率对每个 \(s>0\) 都成立，常数可依赖 \(s\)。
\end{proposition}
\begin{proof}
范数界由两个平均的界直接相加得到。对低频区域 \(|\xi|<R/2\)，令 \(v=|\xi|^2/R^2\)，则组合符号为
\[
 c_\delta(v)=2(1-v/2)^\delta-(1-v)^\delta
 =1-\frac{\delta(\delta-1)}4v^2+O(v^3).
\]
更精确地，\(c_\delta(0)=1\)、\(c_\delta'(0)=0\)，所以 Taylor 积分余项把 \(1-c_\delta(v)\) 写成
\(v^2b(v)\)，其中 \(b\) 在 \([0,1/4]\) 连同所有所需导数有界。

使用\cref{b2:prop:24-lp-rate}证明中的同一低频截断 \(P_R\)，并记 \(g=J^sf\)。高频输入估计
\[
 \|(I-P_R)f\|_p\leq C R^{-s}\cdot\|g\|_p
\]
实际上对任意 \(s>0\) 成立，因为该处仅使用 \(|\xi|\geq R/4\)，不使用 \(s\leq2\)。对它应用 \(I-C_R\)，代价至多为 \(1+3M\)。

低频归一化误差乘子现在是
\[
 R^{s-4}|\xi|^4\cdot b(|\xi|^2/R^2)\cdot
 \chi(\xi/R)(1+4\pi^2|\xi|^2)^{-s/2}.
\]
对 \(|\alpha|\leq d+2\)，其第 \(\alpha\) 阶导数乘以 \(|\xi|^{|\alpha|}\) 后，在 \(0<|\xi|\leq1\) 上由
\(C_\alpha R^{s-4}|\xi|^4\) 控制，在
\(1\leq|\xi|\leq R/2\) 上由
\(C_\alpha(|\xi|/R)^{4-s}\) 控制。理由与前一命题相同：原点的四阶零点控制小频率，归一化链式求导控制大频率。因为 \(0<s\leq4\)、\(R\geq1\)，这些界与 \(R\) 无关。Mihlin 定理给出低频误差的 \(R^{-s}\) 界，再合并高频部并由稠密性推广。

若 \(\delta=1\)，则
\[
 2(1-v/2)-(1-v)=1
\]
在整个低频区域精确成立，低频误差为零；只剩已经对任意 \(s>0\) 成立的高频估计。这证明最后一个断言，但仍使用了一阶平均原先具有统一 \(L^p\) 界的假设。
\end{proof}

组合权有负数，所以它不保留正性，也不能由正权阶数单调公式读出。它改变的是球内部的逼近阶数，而不是每条球边界的曲率障碍。特别是 \(\delta=1\) 时，精确低频抵消改善逼近阶数的作用，仍限于该乘子有界的指数范围。

\ProseParagraphBreak
通常四阶速率也会饱和。取非零 Schwartz 函数 \(f\)，使 \(\widehat f\) 光滑紧支撑，并固定 \(\delta>0\)、\(\delta\ne1\)。当 \(R\) 足够大时，整个频率支撑都在低频区；上面的 Taylor 展开连同频率导数一致成立，所以逆变换在 Schwartz 拓扑中给出
\[
 R^4(C_R-I)f
 \longrightarrow
 -\frac{\delta(\delta-1)}{64\pi^4}(-\Delta)^2f.
\]
系数来自 \(v^2=|\xi|^4/R^4\) 与
\((-\Delta)^2\) 的符号 \(16\pi^4|\xi|^4\)。该极限非零，因此不能只要求输入再光滑一些，就把一般两半径组合的阶数提高到 \(s>4\)。若要继续消去误差，必须增加组合条件或改变平均；稳定的算子界与低频逼近阶数应分别设计。
